%=============================================================
% CONCLUSION
%=============================================================

Searching network architecture, weight precision, and scheduler settings of
a microcontroller as one problem found the accuracy--energy front more
reliably than every alternative tested. Joint qNEHVI reached a median of
$100\%$ of the reference hypervolume in $100$ evaluations and $90\%$ in a
median of $28$, ahead of qNParEGO ($99.6\%$), MOTPE ($91.0\%$), Random
($74.8\%$), and NSGA-II ($56.0\%$) over $20$ seeds. A two-stage search that fixes
the scheduler before the network stalled at $91.2\%$ in all seeds, because the
fixed clock excluded three of the four front designs; joint search beat it by
$8.8$ points.

The evaluation behind these results was checked against silicon where
that was possible. Cycles for eight networks outside the calibration set were
predicted within $7.6\%$, and $4.8\%$ on average, with a bias toward
underestimation of $4.5\%$. Flash and RAM errors were $1.6\%$ and $6.6\%$
on average. Energy is a datasheet model and was not measured, and the int16
and int4 cycle ratios come from the literature.

For a designer, the practical finding is that clock and priority belong in
the same search as the network. Three of the four best designs run at
$120$\,MHz, a clock the Sequential search excluded by fixing $180$\,MHz first.
Where one evaluation costs seconds or more, as a training run or a board
measurement does, a joint model-based search reached $90\%$ of the reference
hypervolume in a median of $28$ evaluations, while MOTPE and
Sequential needed $54$ and $65$, and Random and NSGA-II reached it in one
seed of $20$.

Three limits bound the conclusions. The reference front has four designs
among $378$ acceptable, so the space is small, and the margin over Random
belongs to this problem. Accuracy is one training seed per architecture. The
workload is a synthetic task set on a single core. The next steps are
measuring energy on a board with a current probe, benchmarking int4 and int16
kernels, and repeating the comparison on a larger space, where the two-stage
cap and the margin over Random can be tested again.
